\chapter{Manual de código fuente}\label{manual-codigo}
\begin{refsection}

\section{Introducción}

Este anexo resume la organización del código fuente de Nuvlo. El proyecto se ha implementado como una aplicación web en arquitectura cliente-servidor, con frontend en React, backend en Node.js/Express, persistencia en PostgreSQL mediante Prisma y uso de Redis como caché temporal y soporte operativo.

El código fuente completo se mantiene en el repositorio GitHub del proyecto:

\begin{center}
\url{https://github.com/Bania1/tfg}
\end{center}

De acuerdo con las indicaciones de la plantilla, no se incluye aquí una impresión completa del código, ya que el tribunal puede consultarlo en el repositorio. Este anexo se centra en explicar su estructura, instalación, ejecución, pruebas y principales decisiones técnicas.

\section{Estructura del repositorio}

El proyecto usa un monorepo sencillo con \texttt{npm workspaces}. Esta organización permite evolucionar conjuntamente frontend, backend, paquetes compartidos, documentación, Prisma y scripts.

{\footnotesize
\setlength{\DTbaselineskip}{10pt}
\dirtree{%
.1 tfg.
.2 apps.
.3 web.
.4 src.
.5 components.
.5 views.
.5 lib.
.4 e2e.
.3 api.
.4 src.
.5 services.
.5 security.
.5 cache.
.5 config.
.5 db.
.2 packages.
.3 shared.
.2 prisma.
.3 schema.prisma.
.2 scripts.
.2 docs.
.3 memoria.
.3 decisions.
.2 docker-compose.yml.
.2 package.json.
}
}

\section{Frontend}

El frontend se encuentra en \texttt{apps/web}. Está desarrollado con React y Vite. Las vistas principales se organizan en \texttt{src/views}:

\begin{itemize}
    \item \texttt{Landing.jsx}: pantalla inicial con conexión a Jira y acceso a demo local.
    \item \texttt{DashboardView.jsx}: dashboard principal, filtros, widgets y métricas.
    \item \texttt{BoardView.jsx}: tablero de flujo agrupado por estados.
    \item \texttt{AlertsView.jsx}: reglas de alerta e historial de eventos.
    \item \texttt{ActivityView.jsx}: registro de actividad del sistema.
    \item \texttt{SettingsView.jsx}: conexión con Jira, sesión y proyectos importados.
\end{itemize}

Los componentes reutilizables se encuentran en \texttt{src/components} y la comunicación con el backend se centraliza en \texttt{src/lib/api.js}. La hoja principal de estilos es \texttt{src/styles.css}.

\section{Backend}

El backend se encuentra en \texttt{apps/api}. Está desarrollado con Node.js y Express. Sus responsabilidades principales son:

\begin{itemize}
    \item Gestionar el flujo OAuth 2.0 3LO con Atlassian.
    \item Proteger la sesión mediante cookie \texttt{httpOnly}.
    \item Aplicar protección CSRF en operaciones mutables.
    \item Renovar tokens de Atlassian mediante refresh token rotatorio.
    \item Sincronizar proyectos, issues, sprints y changelog desde Jira Cloud.
    \item Persistir datos en PostgreSQL mediante Prisma.
    \item Calcular métricas y evaluar alertas.
    \item Registrar eventos de actividad y errores relevantes.
\end{itemize}

Los servicios principales se encuentran en \texttt{apps/api/src/services}. Entre ellos destacan \texttt{jiraSync.js}, \texttt{jiraClient.js}, \texttt{projectDashboard.js}, \texttt{metrics.js}, \texttt{alertRepository.js} y \texttt{activityRepository.js}.

\section{Base de datos y Redis}

El modelo de datos se define en \texttt{prisma/schema.prisma}. PostgreSQL almacena información durable: usuarios, sesiones cifradas de Atlassian, proyectos, tableros, sprints, incidencias, transiciones, métricas, alertas, ejecuciones de sincronización y logs de actividad.

Redis se utiliza como componente temporal. En la implementación actual almacena cachés de corta duración para datos de Jira y estados de sincronización. Estos datos tienen TTL y no se consideran fuente principal de verdad. La fuente durable del sistema es PostgreSQL.

\section{Variables de entorno}

La configuración local parte de \texttt{.env.example}. El fichero real \texttt{.env} no debe subirse al repositorio. Las variables más importantes son:

\begin{itemize}
    \item \texttt{DATABASE\_URL}: conexión a PostgreSQL.
    \item \texttt{REDIS\_URL}: conexión a Redis.
    \item \texttt{JWT\_SECRET}: secreto para firmar sesiones.
    \item \texttt{ENCRYPTION\_KEY}: clave para cifrar tokens OAuth.
    \item Credenciales OAuth de Atlassian: identificador de cliente, secreto de cliente y URI de redirección.
    \item \texttt{RETENTION\_*}: parámetros de retención de logs, métricas y eventos.
    \item \texttt{DEMO\_FIXED\_TICK}: opción para congelar la demo offline al generar capturas reproducibles.
\end{itemize}

\section{Instalación y ejecución}

Los comandos básicos de desarrollo son:

\begin{lstlisting}[language=bash]
cp .env.example .env
npm install
docker compose up -d
npm run db:generate
npm run demo:data
npm run db:push
npm run db:seed
npm run dev
\end{lstlisting}

La interfaz web se sirve en \url{http://localhost:5174} y la API en \url{http://localhost:3002}.

\section{Pruebas}

El proyecto incluye pruebas unitarias, pruebas de integración y pruebas E2E:

\begin{lstlisting}[language=bash]
npm test
npm run test:e2e
npm run build
\end{lstlisting}

Las pruebas unitarias validan cálculo de métricas, reglas de alerta, CSRF y política de retención. Las pruebas de integración comprueban endpoints del backend con mocks. Las pruebas E2E con Playwright validan el flujo principal de la UI en modo demo, incluyendo entrada, filtros, widgets y alertas.

\section{Capturas de documentación}

Las capturas de la memoria se generan de forma reproducible mediante scripts:

\begin{lstlisting}[language=bash]
npm run docs:screenshots:jira
npm run docs:screenshots:demo
\end{lstlisting}

El primer comando requiere una sesión OAuth previa y sincroniza el proyecto Jira real antes de capturar. El segundo utiliza mocks y dataset offline para generar imágenes estables de respaldo.

\section{Código de terceros y fuentes externas}

No se han copiado ficheros completos de código fuente externos en la implementación. Se han usado librerías de código abierto instaladas mediante npm, como React, Vite, Express, Prisma, Playwright, Vitest, Recharts y Lucide React, respetando sus licencias y uso habitual como dependencias.

La integración con Jira Cloud se ha implementado consultando la documentación oficial de Atlassian sobre OAuth 2.0 3LO y Jira REST API. Las decisiones técnicas relacionadas con arquitectura, seguridad, pruebas y separación entre Jira real y demo offline se documentan en \texttt{docs/decisions}.

\section{Mantenimiento}

El mantenimiento del proyecto se apoya en commits pequeños, pruebas automatizadas y documentación incremental. Las decisiones arquitectónicas se registran como ADR en \texttt{docs/decisions}. La documentación privada de apoyo sobre cumplimiento, fuentes y uso de IA se mantiene localmente hasta ser revisada con los tutores.

\phantomsection
\addcontentsline{toc}{section}{Referencias}
\printbibliography[heading=subbibliography]
\end{refsection}
